% GENERATED by tools/gen_thesis_tables.py -- do not edit by hand.
% Edit the source CSV and regenerate; see CLAUDE.md §9.
\textbf{Soja} (disponível) & pH do solo & 0,38 & intervalar & 4,5, 5,5, 6,8, 7,5 \\
 & Declive ($^\circ$) & 0,37 & decrescente & 4, 8 \\
 & Teor de argila (\%) & 0,24 & intervalar & 20, 32, 45, 62 \\
 & Temp. do trimestre mais quente ($^\circ$C) & porta & intervalar & 18, 22, 30, 34 \\
 & Precipitação anual (mm) & porta & intervalar & 800, 1100, 1800, 2200 \\
 & Índice de aridez (P/ETP) & porta & intervalar & 0,4, 0,7, 1,3, 1,6 \\
\textbf{Cana-de-açúcar} (disponível) & Teor de argila (\%) & 0,38 & intervalar & 22, 32, 52, 66 \\
 & Teor de água disponível para as plantas & 0,32 & crescente & 22,1, 33,6 \\
 & Declive ($^\circ$) & 0,30 & decrescente & 5, 10 \\
 & Temp. do trimestre mais quente ($^\circ$C) & porta & intervalar & 20, 23, 32, 36 \\
 & Precipitação anual (mm) & porta & intervalar & 900, 1300, 2000, 2600 \\
 & Meses secos (n) & porta & intervalar & 3, 5, 5, 9 \\
\textbf{Outras culturas anuais} (disponível) & Carbono orgânico do solo (g/kg) & 0,54 & crescente & 6, 20 \\
 & pH do solo & 0,16 & intervalar & 4,5, 5,5, 6,8, 7,5 \\
 & Graus-dia de crescimento & 0,14 & intervalar & 2500, 3500, 5500, 7000 \\
 & Declive ($^\circ$) & 0,12 & decrescente & 4, 8 \\
 & CV da precipitação & 0,04 & decrescente & 0,6, 1,3 \\
 & Precipitação anual (mm) & porta & intervalar & 700, 1000, 1700, 2200 \\
\textbf{Piscicultura} (disponível) & Distância à água sazonal (m) & 0,48 & decrescente & 750, 6000 \\
 & Temp. do trimestre mais quente ($^\circ$C) & 0,22 & intervalar & 22, 26, 32, 36 \\
 & Teor de argila (\%) & 0,14 & crescente & 15, 35 \\
 & Declive ($^\circ$) & 0,11 & decrescente & 2,86, 5,71 \\
 & Densidade de drenagem & 0,05 & crescente & 0, 0,09 \\
\textbf{Pecuária} (apenas exclusão) & Umidade do solo (mm) & 0,70 & crescente & 55, 120 \\
 & Declive ($^\circ$) & 0,17 & decrescente & 6, 15 \\
 & Carbono orgânico do solo (g/kg) & 0,13 & crescente & 8, 18 \\
 & Precipitação anual (mm) & porta & intervalar & 700, 1000, 2000, 2600 \\
 & Meses secos (n) & porta & decrescente & 5, 9 \\
\textbf{Conservação} (nenhuma; compensatória) & Distância a área protegida (m) & 0,35 & decrescente & 5000, 40000 \\
 & Carbono em biomassa acima do solo (Mg~C/ha) & 0,17 & crescente & 10, 20 \\
 & Rugosidade do terreno (m) & 0,17 & crescente & 5, 30 \\
 & Índice de umidade topográfica & 0,07 & crescente & 9,44, 11,56 \\
 & Fração de vegetação nativa & 0,07 & crescente & 0,25, 0,85 \\
 & Índice de aridez (P/ETP) $\dagger$ & 0,06 & crescente & 0,95, 1,25 \\
 & Temp. do trimestre mais quente ($^\circ$C) $\dagger$ & 0,05 & decrescente & 23, 31 \\
 & Densidade de drenagem & 0,03 & crescente & 0,01, 0,08 \\
 & Declive ($^\circ$) & 0,02 & crescente & 2,86, 11,31 \\
\textbf{Geração fotovoltaica} (disponível) & Radiação solar (W/m$^2$) & 0,40 & crescente & 186, 204 \\
 & Temp. do trimestre mais quente ($^\circ$C) & 0,18 & decrescente & 24, 31 \\
 & Log do tempo de viagem à cidade & 0,14 & decrescente & 3,5, 5,4 \\
 & Índice de claridade (cobertura de nuvens) & 0,14 & crescente & 0,47, 0,5 \\
 & Declive ($^\circ$) & 0,08 & decrescente & 2,86, 8,53 \\
 & Componente norte da orientação & 0,06 & crescente & -0,3, 0,5 \\
\bottomrule
